\section{La familia de modelos Claude y las prácticas de ingeniería}

\subsection{Sistema de nombres y estrategia de iteración}

Se nombran con formas poéticas, correspondientes a distintas compensaciones entre rendimiento y costo:
\begin{itemize}
    \item \textbf{Haiku}: el más pequeño, el más rápido, el más barato — \textit{``sorprendentemente inteligente''}
    \item \textbf{Sonnet}: intermedio — más inteligente pero más lento y más caro
    \item \textbf{Opus}: el más grande, el más potente — la versión más inteligente al momento de su lanzamiento
\end{itemize}

La estrategia central es \textbf{desplazar la curva de compensación}: cada nueva generación no solo mejora en un punto, sino que desplaza toda la curva hacia arriba y a la derecha. Sonnet 3.5 ya es más inteligente que el Opus 3 original (sobre todo en código), y Haiku 3.5 es más o menos equivalente a Opus 3. \textit{``El objetivo es desplazar esa curva.''}

El nombramiento en sí mismo resultó un desafío inesperado. Los modelos de IA no son como el software tradicional (versiones incrementales 3.7, 3.8), porque las mejoras del preentrenamiento pueden llegar más rápido de lo previsto y alterar el plan de nombres. \textit{``Nadie ha resuelto el nombramiento... este es, de alguna manera, un paradigma distinto del software común.''}

\subsection{El trabajo entre un lanzamiento de modelo y otro}

Entre cada lanzamiento hay una gran cantidad de trabajo invisible:

\begin{enumerate}
    \item \textbf{Preentrenamiento}: meses, con decenas de miles de GPU o TPU
    \item \textbf{Entrenamiento posterior}: RLHF y otros métodos de RL, \textit{``a una escala cada vez mayor''}
    \item \textbf{Pruebas de seguridad}: pruebas internas de RSP más instituciones externas (los institutos de seguridad de IA de Estados Unidos y el Reino Unido, pruebas de terceros sobre riesgos CBRN)
    \item \textbf{Optimización de inferencia}: hacer que el modelo funcione de manera eficiente en producción
    \item \textbf{Puesta en marcha de la API}: despliegue y lanzamiento
\end{enumerate}

\textit{``Te sorprendería cuántos de los retos se reducen a ingeniería de software y a ingeniería de rendimiento... casi siempre se reduce a los detalles, y a menudo a detalles muy, muy aburridos.''} No es un momento repentino de inspiración tipo Eureka.

\subsection{``¿Claude se volvió tonto?''—— el efecto psicológico del Wi-Fi}

Esta es una de las quejas más comunes en Reddit. Dario lo aclara sin ambigüedad: sin el anuncio de un nuevo modelo, \textbf{los pesos del modelo no cambian}.

\begin{warningbox}{Por qué todo mundo siente que el modelo se volvió tonto}
Este es un \textbf{fenómeno generalizado} — todas las empresas de modelos de frontera (GPT-4, Claude, Gemini) enfrentan la misma queja. Explicaciones posibles:

\begin{itemize}
    \item \textbf{Efecto psicológico}: una vez que se desvanece la novedad, los defectos se vuelven más notorios
    \item \textbf{Sensibilidad a la redacción}: un cambio mínimo en cómo se formula algo puede producir resultados distintos
    \item \textbf{Cambios en el System Prompt}: activar por defecto la función Artifacts modificó el System Prompt
    \item \textbf{Memoria selectiva}: se recuerdan las malas experiencias y se olvidan las buenas
\end{itemize}

Dario hizo una analogía muy acertada con el Wi-Fi de los aviones: \textit{``La primera vez que hubo Wi-Fi en un avión, se sentía como magia; ahora es 'esta porquería otra vez no funciona, qué basura'.''} Cuando la línea base sube, las expectativas suben con ella.

Las pruebas AB solo se ejecutan brevemente antes y después del lanzamiento de un modelo. Dario reconoce que el System Prompt cambia de vez en cuando, pero los pesos en sí no cambian.
\end{warningbox}

\subsection{El acertijo del ``topo que golpear'' en la personalidad de Claude}

Gestionar el comportamiento de Claude es un problema de optimización multidimensional: corregir un problema suele provocar otro.

\begin{knowledgebox}{La dificultad intrínseca de controlar el comportamiento}
Corregir la verbosidad $\to$ el modelo se vuelve perezoso (escribe ``el resto del código va aquí'' en el código). Reducir los rechazos $\to$ aumentan los problemas de seguridad.

\textit{``No es que queramos ahorrar cómputo... es que controlar el comportamiento de un modelo es realmente difícil.''} Esta es \textit{``una analogía en el presente del problema del control de la IA del futuro''} —— si hoy ya nos cuesta trabajo lograr que un modelo trace la línea entre ``ayudarte con tu clase de posgrado en virología'' y ``no ayudarte a fabricar viruela'', ¿qué tan difícil será controlar una IA más poderosa en el futuro?

\textit{``Le pides al modelo que no ayude a la gente a fabricar y propagar la viruela, pero que sí esté dispuesto a ayudarte en tu clase de posgrado en virología —— ¿cómo lograr las dos cosas a la vez? Es muy difícil.''}
\end{knowledgebox}

\subsection{Resumen de la sección}

La familia Claude itera desplazando de manera constante la curva de compensación entre rendimiento y costo. Que se sienta ``más tonto'' es un efecto psicológico, no un deterioro real. El acertijo del ``topo que golpear'' en el comportamiento del modelo es un reflejo, en el presente, del problema del control de la IA.
